

\subsection*{Author inventory alternatives on the same annual book}

The optional annual inventory method can now be described in source text and
compared across named scenarios.  These interfaces extend the checked execution
above.  They retain its opening balances, tax liabilities and recognition
assumptions, so a change in how a sale is written cannot silently change the
company being simulated.

\texttt{Annual\allowbreak{}Inventory\allowbreak{}Authoring/1} uses the annual declaration grammar with an
explicit \texttt{partial\_\allowbreak{}inventory\_\allowbreak{}sale} operation.  Its header names the sale;
the body supplies the sale profile, date, allocation method, lot, units sold and
cash proceeds.  Each of the six inventory findings remains a separate
known-or-unknown declaration.  A known finding carries the text explaining its
provenance; an unknown finding carries a reason.  This inventory convention is
preserved alongside the structured provenance used by the older tax
operations.  Neither form is silently translated into the other.

Compilation binds the declared opening request to its exact content digest and
returns the same \texttt{Annual\allowbreak{}Inventory\allowbreak{}Program/1} request as the typed form.
It does not execute a sale, assess tax or advance a period.  Consequently, a
false divisibility finding or an unresolved homogeneity finding survives as
explicit data until the financial evaluator refuses the operation.  Missing
fields, duplicate declarations and malformed rational quantities fail earlier,
when the source is read.  The compiled record retains the original text hash
and the lines containing each operation and field.  Whitespace can change
that text hash without changing the compiled financial request.

\texttt{Annual\allowbreak{}Inventory\allowbreak{}Family/1} then gives each alternative its own name,
assumptions and ordered operations, all starting from the same declared annual
book.  Consider the earlier example with two recovered units costing 20 and
cash of 452.  The first alternative sells one unit for 15, assesses current tax
of 2, refreshes deferred measurement and carries the remaining unit into the
following year.  Its first-year profit is 3 and closing cash is 467.  Selling the
last unit for 15 in the next year gives that year's profit of 5 and closing cash
of 482, while the carried tax payable remains \(49/4\).  The two annual statement
vectors belong to this one history and are returned together with its name.

A different alternative may fail before the carry because the homogeneous-lot
finding is unknown, or may reach the second year and attempt to sell more units
than remain.  The family retains the precise valid prefix and the operations
that were not executed.  A report for a year that was never reached is
unavailable; it is not a statement with zero balances.  An explicit budget can
also leave members unvisited, while preserving their names and assumptions.
Even a zero member budget checks the common opening request.  Equal financial
results do not merge separately named alternatives.

Joint reports recompute the saved family before returning the complete archived
or current-prefix statements.  The reference checks cover inventory quantity,
current-tax histories, residual-stock carry, source locations, rejected premises
and retained alternatives.  These are authored, unweighted scenarios under the
selected proportional inventory model.  They establish neither probabilities
for those scenarios nor membership of an actual company's unknown history.
General inventory policies and independent accounting acceptance remain open.
